% Imporditud: tools/import_eksperimendid.py
% Allikas: Eesti füüsikaolümpiaadi eksperimendi ülesannete
%   kogu (Rael Kalda, 2023), ülesanne Ü123, lahendus L123.
\setAuthor{Erkki Tempel}
\setRound{lõppvoor}
\setYear{2021}
\setNumber{P E2}
\setDifficulty{3}
\setTopic{varia}
\setVahendid{Anneke šokolaad 300g, käärid}
\setPunktid{12}
\setAllikas{Eksperimendi kogu Ü123, lahendus lk 97}
\setLahendusseis{toores}

\prob{Paberi tihedus}
Mitu korda erinevad šokolaadi ümbrispaberi ja hõbepaberi tihedused?

\hint

\solu
Lõikame mõlemad paberid täpselt sama suurusteks, mida suurem, seda parem. Tiheduste suhte leidmiseks peame leidma fooliumpaberi ja ümbrispaberi masside suhte mf /mp ning paksuste suhte hf /hp. Tiheduse valemist saame avaldada tiheduste suhte masside ja paksuste suhete kaudu

$\rho$f = mf Vp = mf $\cdot$ hp $\rho$p mpVf mp hf

Masside suhte leidmiseks kasutame kangimeetodid. Paberi saame voltida nii, et seda saab kangina kasutada ja saame leida paberite masside suhte mf /mp. Paberi paksuste suhte hp/hf leidmiseks voldime paberi mitu korda kokku (näiteks 16 kihti) ning vaatame, mitu kihti paberit on sama paks kui 16 kihti fooliumpaberit.
\probend
